\chapter{双稳态、开关与记忆}
现在让产生率依赖蛋白质自身。单个动态变量只能沿浓度轴运动，但产生曲线的形状已经能决定系统是回到唯一表达水平，还是在不同长期状态之间作出选择。

\begin{figure}[H]\centering
\begin{tikzpicture}
\node[gene] (none) at (0,0) {$X$};\node at (0,-.8) {无反馈};
\node[gene] (pos) at (4,0) {$X$};\node at (4,-.8) {自激活};
\node[gene] (neg) at (8,0) {$X$};\node at (8,-.8) {自抑制};
\draw[act] (pos) to[loop above,min distance=12mm] (pos);
\draw[rep] (neg) to[loop above,min distance=12mm] (neg);
\end{tikzpicture}
\caption{单节点的三种符号拓扑。所有模型仍包含产生与降解；自边表示自身浓度对产生率的作用。}
\end{figure}
\section{单基因反馈：恢复、双稳态与迟滞}
无反馈模型的产生率恒为 $\alpha$。自抑制把它改成随 $x$ 非增的 $F(x)$，例如 $F(x)=a/(1+(x/K)^n)$。浓度上升时，产生减少而降解增加，两种作用都将浓度拉回。

把这种直觉写成方程，令 $f(x)=F(x)-\gamma x$，假设产生率 $F(x)$ 非负、有界、非增且为 $C^1$。这里“有界”就是有一个有限上限，“非增”表示浓度增加时产生率不会上升，$C^1$ 则保证可以连续地计算斜率。$f$ 就严格递减，且 $f(0)\geq0$、$f(x)\to-\infty$。因此产生曲线与降解直线恰好相交一次，稳态 $x^*$ 唯一，允许它位于边界 $x^*=0$。

不仅稳态唯一，所有非负初值也都会趋于它。令浓度偏差 $e=x-x^*$，则
\[
\frac12\frac{\dd}{\dd t}e^2=e\{F(x)-F(x^*)\}-\gamma e^2\leq-\gamma e^2,
\]
因为 $F$ 非增，浓度偏差与产生率偏差的乘积不可能为正。积分可得
\[
|x(t)-x^*|\leq e^{-\gamma t}|x(0)-x^*|.
\]
式~\eqref{eq:bound}说明浓度不会无限增大，所以这个恢复速度的估计可以一直使用。若 $F'(x^*)<0$，小扰动满足 $\dot e\simeq(F'(x^*)-\gamma)e$，小偏差减小到原来约 $37\%$ 所需的时间为 $1/(\gamma-F'(x^*))$，短于无反馈的 $1/\gamma$。这把“自抑制促进恢复”落实为具体速率；它描述的是稳态附近的小偏差；输入变化很大时，还要直接计算整条时间轨迹。


自激活的方向相反：浓度上升也会提高产生率，偏差可能被放大。于是一个自然的问题是，这种放大能否使系统保留不同表达状态？取
\begin{equation}\dot x=a+b\frac{x^n}{K^n+x^n}-\gamma x.\label{eq:auto}\end{equation}
稳态是产生曲线与线性降解直线的交点。局部稳定条件为
\begin{equation}F'(x^*)-\gamma<0.\end{equation}
图~\ref{fig:auto} 的参数 $a=0.05,b=3.5,K=\gamma=1,n=4$ 产生三个稳态，外侧稳定，中间不稳定。关键不是正反馈必然双稳态，而是正反馈与足够的非线性产生多个交点。改变参数时，两个交点可能靠近并合在一起，再一起消失。合并的临界位置称为“折叠点”：此时产生曲线与降解直线不仅相交，斜率也相同，即 $F(x^*)=\gamma x^*$ 和 $F'(x^*)=\gamma$。
\begin{figure}[H]\centering\includegraphics[width=.74\linewidth]{figures/self_activation.pdf}\caption{自激活模型的产生率与降解率，实心点为稳定的稳态，空心点为不稳定的稳态。}\label{fig:auto}\end{figure}

先判断交点两侧的运动方向，再问交点何时会成对出现。对只有一个浓度变量的 $C^1$ 方程，若 $f'(x^*)<0$，则足够靠近稳态时 $f$ 在左侧为正、右侧为负，轨道向稳态运动；若 $f'(x^*)>0$ 则方向相反。如果交点处的斜率为零，就需要再看它左右两侧的实际变化方向，不能只凭这个斜率判断。对式~\eqref{eq:auto}，
\[
F'(x)=\frac{bnK^nx^{n-1}}{(K^n+x^n)^2}.
\]
当 $n>1$ 时该导数的唯一最大值出现在 $x_m=K((n-1)/(n+1))^{1/n}$：令 $h=x/K$，对 $h^{n-1}/(1+h^n)^2$ 求导，驻点条件为 $(n-1)-(n+1)h^n=0$。因此
\[
\max F'=\frac bK\frac{(n+1)^2}{4n}
\left(\frac{n-1}{n+1}\right)^{(n-1)/n}.
\]
若该值小于 $\gamma$，$f$ 严格递减，稳态唯一；超过 $\gamma$ 只是出现多交点的必要条件，仍需选择合适的 $a$。
令 $q=(x^*/K)^n$，联立折叠条件可参数化为
\[
b_{\rm fold}=\frac{\gamma K}{n}q^{(1-n)/n}(1+q)^2,\qquad
a_{\rm fold}=\frac{\gamma K}{n}q^{1/n}(n-1-q).
\]
这是由 $F'=\gamma$ 先解出 $b$，再代入 $a=\gamma x-bq/(1+q)$ 得到。对 $a>0$ 需 $0<q<n-1$。通常的折叠分岔还要求 $f''\ne0$，称为一个\emph{非退化条件}。“分岔”指参数变化使稳态或轨迹的类型发生改变；“非退化”在这里指相切处仍有明确的弯曲，两个交点的合并不是更特殊的平坦情形。若连弯曲程度也为零，交点怎样变化就需要进一步分析。

把条件落实到图中的参数，也能解析地确认这三个交点。取 $f(x)=0.05+3.5x^4/(1+x^4)-x$，有 $f(0)>0$、$f(0.1)<0$、$f(1)>0$、$f(4)<0$，所以三个互不相交的区间 $(0,0.1)$、$(0.1,1)$、$(1,4)$ 各含至少一个根。上面的计算说明 $f'$ 最多有两个零点，也就是说曲线最多改变两次升降方向。若 $f$ 有四个交点，相邻交点之间就至少需要三个斜率为零的位置，超过了这个上限。这用到了\emph{Rolle 定理}：一条可求导的曲线在两端取相同值时，中间至少有一处斜率为零。因此三个交点就是全部结果。两个极值分别严格低于、严格高于零，因此三个根均不在极值点处，斜率依次为负、正、负。中间的不稳定交点构成一道分界：初始浓度在它以下时趋于低表达态，在它以上时趋于高表达态。短暂扰动跨过此阈值可造成持久状态切换。


上述例子有两个稳定交点，因此具有\emph{双稳态}。这里的“双”指同一组参数、同一个固定外部输入下，有两个能吸引附近轨迹的长期状态；两个交点本身还不够，必须检查它们的稳定性。下面把“可能出现三个交点”进一步变成可直接判断的参数条件。

固定 $b,K,n,\gamma$，假设 $n>1$ 且 $\max F'>\gamma$。因为 $F'$ 先增后减，方程 $F'(x)=\gamma$ 有两个正根 $x_-<x_+$。于是 $f=F-\gamma x$ 在 $x_-$ 处取局部最小值，在 $x_+$ 处取局部最大值。令
\[
a_{\rm up}=\gamma x_--b\frac{x_-^n}{K^n+x_-^n},\qquad
a_{\rm down}=\gamma x_+-b\frac{x_+^n}{K^n+x_+^n}.
\]
则 $f(x_-)=a-a_{\rm up}$、$f(x_+)=a-a_{\rm down}$。两个极值之间 $f$ 严格增加，所以 $a_{\rm down}<a_{\rm up}$。在物理范围 $a>0$ 内，\emph{恰有三个不同稳态的充要条件}是
\[
a_{\rm down}<a<a_{\rm up}.
\]
证明可以按三个单调区间逐段看：$f(0)=a>0$，第一个极小值为负，随后极大值为正，最后 $f(x)\to-\infty$，每一段恰好经过零一次。反过来，若有三个不同交点，极小值必须低于零、极大值必须高于零，否则至少一段没有交点。这就补上了“正反馈足够强”之后仍需检查的截距条件。若下阈值为负，非负的基础产生率未必能跨过它；实验上能否完整观察迟滞，还要看实际可调参数的范围。

记三个根为 $x_L<x_U<x_H$。变化率的符号依次是
\[
\begin{array}{c|cccc}
x &(0,x_L)&(x_L,x_U)&(x_U,x_H)&(x_H,\infty)\\\hline
\dot x&+&-&+&-
\end{array}
\]
因此 $x_L,x_H$ 渐近稳定，$x_U$ 不稳定。进一步地，唯一性保证轨迹不能穿过一个稳态；位于 $x_U$ 以下的轨迹向 $x_L$ 单调运动，位于其以上的轨迹向 $x_H$ 单调运动。单调且有界的轨迹有极限；若极限处 $f$ 不为零，变化率就会在附近保持一个非零大小，轨迹不可能停下来，所以极限必为相应稳态。由此得到完整的\emph{吸引域}：$[0,x_U)$ 属于低态，$(x_U,\infty)$ 属于高态，而初值恰为 $x_U$ 才会一直停在分界上。吸引域就是“哪些初始状态最终进入同一个长期状态”的集合。

还可以定义一维\emph{势函数}
\[
V(x)=-\int_0^x f(\xi)\dd\xi,\qquad
\frac{\dd V(x(t))}{\dd t}=-f(x(t))^2\leq0.
\]
稳定交点对应势谷，不稳定交点对应势垒。这个图像解释了为何小扰动会被拉回、跨过阈值的扰动却能改变长期状态。这里的 $V$ 是从方程构造的数学函数，通常不是分子的热力学自由能；两个势谷也不要求等深，才能形成确定性的双稳态。没有额外扰动时，轨迹不会自行跨过分界。

\emph{迟滞}进一步涉及改变参数，而不仅是选择不同初值。若缓慢提高基础产生率 $a$，低态会沿自身分支保持到 $a_{\rm up}$，再跳到高态；反向降低 $a$，高态则保持到 $a_{\rm down}$ 才消失。两次切换阈值不同，所以相同输入下的输出取决于此前是从低态还是高态到来。也可以固定 $a$、调节反馈强度 $b$，得到图~\ref{fig:bistability_branches}。参数变化必须相对恢复足够慢，才能近似跟随稳态分支；较快扫描还可能带来额外的动态滞后，不能把所有滞后都解释为静态双稳态。
\begin{figure}[H]\centering
\includegraphics[width=\linewidth]{figures/bistability_branches.pdf}
\caption{把双稳态、阈值与迟滞放在同一张图中。左：固定 $a=0.3,K=\gamma=1,n=4$，随反馈强度 $b$ 改变的稳态分支；实线稳定、虚线不稳定，箭头表示慢扫描时的切换方向。中：双稳态区中的运动方向及吸引域分界。右：相同参数下构造的势函数，两个谷之间由势垒隔开。数值只用于绘图，分支和吸引域由正文论证。}\label{fig:bistability_branches}
\end{figure}

折叠点附近还会出现\emph{临界减慢}。设控制参数为 $a$，折叠点为 $(x_c,a_c)$，此时 $f=f_x=0$、$f_{xx}\ne0$ 且 $f_a=1$。令 $\eta=x-x_c$、$\mu=a-a_c$，Taylor 展开给出
\[
\dot\eta=\mu+\frac12f_{xx}(x_c,a_c)\eta^2
 +O(|\eta|^3).
\]
“$O$”表示被省略的项至多按括号内的量级增长。稳态靠近临界点时，两个分支的距离按 $\sqrt{|\mu|}$ 缩小，稳定分支上的恢复斜率 $|f_x|$ 也按 $\sqrt{|\mu|}$ 趋于零，故小扰动恢复时间约为 $1/|f_x|$，会越来越长。这个结论说明：开关尚未切换之前，就可能变得恢复很慢。它来自光滑、非退化折叠的局部展开；退化点需保留更高阶项。

两个稳定表达态已经足以构成记忆，却还不足以构成时钟。对\emph{一维自治方程} $\dot x=f(x)$，“一维”表示只有一个动态变量，“自治”表示变化率由当前状态决定、不显式依赖时间。如果外部条件固定、也没有延迟，那么同一个浓度 $x$ 总对应同一个变化方向。浓度在一个区间里一直上升，就不能在没有其他变化的情况下突然掉头下降。

可以把这一点完整地论证出来。假设有一条重复往返的轨迹，周期为 $T$。它在一个周期内一定有最高浓度 $x_m$。最高点处的变化率为零，所以 $f(x_m)=0$，也就是说这里已经是稳态。从这个浓度出发，始终停在 $x_m$ 就满足原方程；同一初始浓度又只能产生一条轨迹，因此原来的轨迹也只能停在那里，与“重复往返”矛盾。

所以，一个变量可以提供不同表达状态之间的选择，但在这些假设下不能自行维持振荡。要寻找持续振荡，需要增加动态环节，或者另行考虑延迟、随时间变化的输入等因素。
\section{互抑开关：对称破缺与状态选择}
负反馈倾向于恢复，正反馈则可能放大两个节点之间的差异。互抑制尤其直观：$X$ 增多会压低 $Y$，较低的 $Y$ 又解除对 $X$ 的抑制，最初的表达优势于是被增强。用下面的无量纲模型分析这一过程：
\begin{equation}\dot x=\frac{a}{1+y^n}-x,\qquad\dot y=\frac{a}{1+x^n}-y.\label{eq:toggle}\end{equation}
先把两条变化率分别设为零，得到 $x=a/(1+y^n)$ 与 $y=a/(1+x^n)$。这样的曲线称为“零流线”：在一条线上 $x$ 暂时不变，在另一条线上 $y$ 暂时不变；两条线的交点处两个变量都不变，就是稳态。图~\ref{fig:toggle} 展示 $a=3,n=3$ 的示例：两个不对称表达态和中间的对称稳态。再检查这些交点附近的小偏差会恢复还是被放大，就能判断它们是否稳定。该模型展示状态选择的最小机制，但生物学中的细胞命运并不等同于某一对基因浓度。

\begin{figure}[htbp]\centering\includegraphics[width=.67\linewidth]{figures/toggle_nullclines.pdf}\caption{互抑开关的零流线与相流（$a=3,n=3$）。实心点稳定，空心点为鞍点；绿色轨迹分别趋于两种不对称表达态。}\label{fig:toggle}\end{figure}

先从两个基因表达相同的状态开始，问微小的不对称是否会增长。令 $x=y=s$，则 $a=s(1+s^n)$；右侧在 $s\geq0$ 严格递增，故对称稳态唯一。记交叉斜率的绝对值为
$h=ns^n/(1+s^n)$，雅可比为 $\begin{psmallmatrix}-1&-h\\-h&-1\end{psmallmatrix}$。
如果两个浓度同时增加同样的一小点，偏差按 $e^{(-1-h)t}$ 衰减；如果一个增加、另一个减少同样的一小点，差异则按 $e^{(-1+h)t}$ 变化。这分别对应特征值 $-1-h$ 与 $-1+h$。当 $n>1$ 时，后者过零的阈值为
\[
s_c=(n-1)^{-1/n},\qquad a_c=\frac n{n-1}(n-1)^{-1/n}.
\]
当 $n\leq1$ 时有限 $s>0$ 下 $h<1$，对称态不会经这一线性机制失稳。阈值只判断对称态；不对称态的存在和稳定性仍须证明。

对于图中的 $a=3,n=3$，可以继续把所有稳态解出来，从而确认系统有一个对称鞍点和两个稳定的不对称表达态。
稳态满足 $x+xy^3=3$、$y+yx^3=3$，故 $x,y>0$。相减得
$(x-y)[1-xy(x+y)]=0$。对称解满足 $s+s^4=3$，严格单调性保证唯一；又 $s>1$，故 $h=3s^3/(1+s^3)>1$，是鞍点。
不对称时令 $S=x+y$、$P=xy$，有 $PS=1$。相加得
$S+P(S^2-2P)=6$，即 $S-S^{-2}=3$。左侧在 $S>0$ 严格递增，故恰有一个解，且 $S>3$。二次方程 $t^2-St+1/S=0$ 的判别式 $S^2-4/S>0$，两根均正；交换两根得恰好两个不对称解。此和、差条件亦反推出原两条稳态方程，未引入额外解。
在这些点，两个交叉斜率均等于 $-3x^2y^2/a=-P^2$，特征值为 $-1\pm P^2$。因为 $P=1/S<1/3$，两根均负。三种分支已穷尽所有稳态。小的不对称在鞍点附近被放大，而两个不对称态附近的偏差又被压回去；这两步共同构成了开关的状态选择与保持。

只计算 $a=3$ 还看不出开关怎样出现。对于整个 $n=3$ 参数族，上述消元可以原样保留：不对称稳态必须满足
\[
P=1/S,\qquad a=S-S^{-2},\qquad
x,y=\frac{S\pm\sqrt{S^2-4/S}}2.
\]
函数 $S-S^{-2}$ 在 $S>0$ 严格递增，而两个不同正根存在恰好要求 $S^3>4$。因此不对称稳态出现的阈值为
\[
S_c=\sqrt[3]4,\qquad a_c=S_c-S_c^{-2}=\frac3{2^{4/3}},
\qquad x=y=s_c=2^{-1/3}.
\]
$a<a_c$ 时只有对称稳态，且 $h<1$、两根特征值都为负；$a>a_c$ 时对称态变成鞍点，并且恰好出现两个不对称稳态。后两者的交叉斜率相等，绝对值为 $3P^2/a=3/(S^3-1)$，故
\[
\lambda_\pm=-1\pm\frac3{S^3-1}<0\qquad(S^3>4).
\]
这同时证明了两条新分支的存在与局部渐近稳定，而非只从对称态失稳猜测它们存在。

这种从对称状态分出两种对等选择的变化称为\emph{超临界叉形分岔}：两条稳定分支在临界点连续长出。写表达差 $D=x-y$，有 $D^2=S^2-4/S$。在临界点处
\[
\frac{\dd a}{\dd S}=\frac32,\qquad
\frac{\dd(D^2)}{\dd S}=3S_c,
\qquad
D=\pm\sqrt{2S_c(a-a_c)}+O((a-a_c)^{3/2}).
\]
正负两号对应哪一个基因取得表达优势。$a=a_c$ 处线性化有零特征值，所以不能仅用线性斜率下结论。令 $m=(x+y)/2$、$d=(x-y)/2$，把快速恢复的共同表达变量写成 $m=s_c+Md^2+O(d^4)$。在临界点，$f'(s_c)=-1$、$f''(s_c)=0$，故共同表达方程的二阶项给出 $M=0$；直接求导又有 $f'''(s_c)=8s_c$。差异方程因此化为
\[
\dot d=-\frac{4s_c}{3}d^3+O(d^5).
\]
这说明足够小的差异仍被压回零，但恢复是缓慢的代数过程。在领先近似下，
\[d(t)=\frac{d_0}{\sqrt{1+\frac{8s_c}{3}d_0^2t}}.\]
恢复不再呈指数形式。这里的\emph{中心流形}近似将快变量随慢差异一起消去，所需系数已在上述展开中算出。越过临界点后，线性增长与三次恢复竞争，形成两条稳定分支。

开关的分界在这个对称模型中也能精确定位。交换 $x,y$ 不改变方程，直线 $x=y$ 正向不变，在其上只有一维恢复到对称稳态。对 $a>a_c$，鞍点的稳定方向为 $(1,1)$，不稳定方向为 $(1,-1)$；局部稳定流形恰好就是这条对角线。又令 $D=x-y$，有
\[
\dot D=D\left[\frac{a(x^2+xy+y^2)}{(1+x^3)(1+y^3)}-1\right].
\]
因此 $D$ 的符号保持不变。轨迹有界且最终远离坐标轴，散度恒为 $-2$，排除了周期轨道。两个不对称稳态是吸引点，没有出向连接；鞍点的入向轨迹只在对角线上，也不存在闭合的稳态连接。Poincar\'e--Bendixson 定理将此类有界平面轨迹的长期极限限制为稳态、周期轨道或稳态连接（参考文献~7）；后两类已排除。鞍点仅吸引对角线上的轨迹，故 $x>y$ 的轨迹趋于 $X$ 高表达态，$x<y$ 的轨迹趋于 $Y$ 高表达态。这就证明了两个吸引域及其分界。参数不对称时，对角线通常不再是分界，叉形也会展开成偏向一方的分支，需重新分析。


